\chapter{Constrained Systems and Hamiltonian Motion}

\lectureoverview{Two constrained systems show why Lagrangian mechanics is so
effective: once suitable coordinates have been chosen, the equations of motion
and their conservation laws follow without an explicit calculation of the
constraint forces.  The same examples then lead to Hamiltonian mechanics, in
which coordinates and momenta form phase space, the harmonic oscillator becomes
a rigid rotation, and Hamilton's two first-order equations replace each
second-order Euler--Lagrange equation.}

\section{What It Means to Solve a Mechanical Problem}

The immediate purpose of the Lagrangian method is practical.  A wedge sliding
on a floor and a double pendulum can both be treated with Newton's laws, but the
unknown contact forces and tensions quickly make the bookkeeping unpleasant.
Once the rules of Lagrangian mechanics are known, the same problems become a
sequence of controlled steps:
\begin{enumerate}
  \item choose independent coordinates that label an instantaneous
  configuration;
  \item express the kinetic and potential energies in those coordinates;
  \item form \(L=T-V\) and apply the Euler--Lagrange equations.
\end{enumerate}

For the present purpose, a problem is already \emph{solved} when its equations
of motion have been written down correctly.  One may still want a closed-form
trajectory, but that is a separate question.  Nonlinear differential equations
can be integrated numerically once the mechanical setup is complete.  The hard
part is often not the integration; it is finding variables that incorporate the
constraints without introducing redundant unknown forces.

\section{A Particle Sliding on a Moving Wedge}

Consider a wedge of mass \(M\) constrained to slide without friction along a
horizontal plane.  A point particle of mass \(m\) is constrained to remain on
the inclined face.  There are two degrees of freedom: one tells us where the
wedge is, and one tells us where the particle is on the wedge.

Let
\begin{equation}
  \begin{aligned}
    X&=\text{horizontal position of the wedge corner},\\
    x&=\text{horizontal displacement of the particle}\\[-0.2em]
     &\phantom{={}}\text{from that corner}.
  \end{aligned}
  \label{eq:l6-wedge-coordinates}
\end{equation}
The particle's absolute horizontal coordinate is therefore \(X+x\).  The
choice is not unique.  We could instead use distance along the slope or the
particle's absolute horizontal coordinate.  Any nonredundant pair that labels
the same configurations is legitimate, though some pairs expose the structure
more clearly than others.

For a fixed slope angle \(\beta\), the vertical displacement \(y\) is not an
independent coordinate:
\begin{equation}
  y=x\tan\beta,
  \qquad
  \dot y=\dot x\tan\beta.
  \label{eq:l6-wedge-constraint}
\end{equation}
We specialize to \(\beta=45^\circ\), so that \(y=x\) and
\(\dot y=\dot x\).  This simplification removes trigonometric factors without
changing the point of the example.

\begin{figure}[htbp]
  \centering
  \includegraphics[width=0.84\textwidth]{lecture_06_figure_06.png}
  \caption{The moving wedge, the constrained particle, and the horizontal and
  vertical velocity components used in the \(45^\circ\) example.}
  \lectureframe{00:10:25}
\end{figure}

The wedge has kinetic energy
\begin{equation}
  T_{M}=\frac{M}{2}\dot X^2.
  \label{eq:l6-wedge-kinetic}
\end{equation}
The particle's horizontal velocity is \(\dot X+\dot x\), while its vertical
velocity is \(\dot x\).  Its kinetic energy is consequently
\begin{equation}
  T_m=\frac{m}{2}\left[(\dot X+\dot x)^2+\dot x^2\right].
  \label{eq:l6-particle-kinetic}
\end{equation}
Taking the zero of gravitational potential at \(x=0\), we have \(V=mgx\).
Thus
\begin{equation}
  L=\frac{M}{2}\dot X^2
    +\frac{m}{2}\left[(\dot X+\dot x)^2+\dot x^2\right]
    -mgx.
  \label{eq:l6-wedge-lagrangian}
\end{equation}

\begin{figure}[htbp]
  \centering
  \includegraphics[width=0.86\textwidth]{lecture_06_figure_07.png}
  \caption{The wedge geometry, \(V=mgx\), and the two kinetic-energy
  contributions assembled on the board.}
  \lectureframe{00:14:15}
\end{figure}

\subsection{Canonical Momentum Is Defined by the Lagrangian}

The canonical momenta are defined by the Lagrangian, not by a guess based on a
single particle velocity:
\begin{align}
  p_X&=\frac{\partial L}{\partial\dot X}
      =M\dot X+m(\dot X+\dot x),
      \label{eq:l6-wedge-pX}\\
  p_x&=\frac{\partial L}{\partial\dot x}
      =m(\dot X+\dot x)+m\dot x.
      \label{eq:l6-wedge-px}
\end{align}
The first expression is the ordinary total horizontal momentum:
the wedge contributes \(M\dot X\), and the particle contributes
\(m(\dot X+\dot x)\).  It appears here as the momentum conjugate to the
coordinate \(X\).

\begin{figure}[htbp]
  \centering
  \includegraphics[width=0.86\textwidth]{lecture_06_figure_08.png}
  \caption{The canonical momenta conjugate to the wedge coordinate \(X\) and
  the relative particle coordinate \(x\).}
  \lectureframe{00:18:05}
\end{figure}

The Lagrangian has no explicit dependence on \(X\).  A horizontal translation
of the whole apparatus acts as
\begin{equation}
  \delta X=\epsilon,
  \qquad
  \delta x=0.
  \label{eq:l6-wedge-translation}
\end{equation}
The general infinitesimal symmetry rule
\begin{equation}
  \delta q_i=\epsilon f_i(q),
  \qquad
  Q=\sum_i p_i f_i
  \label{eq:l6-noether-rule}
\end{equation}
therefore gives \(Q=p_X\).  Equivalently, \(X\) is cyclic and
\begin{equation}
  \dot p_X=\frac{\partial L}{\partial X}=0.
  \label{eq:l6-wedge-conservation}
\end{equation}

There is no corresponding vertical translation symmetry.  The wedge is not
allowed to move vertically, and changing the particle's height changes the
gravitational potential.  Horizontal momentum is conserved; vertical momentum
of this constrained subsystem is not.

For the relative coordinate, the Euler--Lagrange equation is
\begin{equation}
  \dot p_x=\frac{\partial L}{\partial x}=-mg.
  \label{eq:l6-wedge-force}
\end{equation}
This line required a sign correction at the board.  The dependable route is
to begin with \(L=T-V\): since \(V=mgx\), the derivative of the Lagrangian is
\(-mg\).  The fact that \(x\) was introduced as a horizontal displacement does
not remove gravity from its equation, because the constraint ties that
horizontal displacement to the vertical height \(y=x\).

No normal force or reaction force has appeared.  Such forces are whatever
they must be to maintain the constraints.  Newton's method would introduce
them and solve for them along with the accelerations; the generalized
coordinates have made them unnecessary.

\section{The Double Pendulum and the Art of Coordinates}

Now suspend two unit masses from two massless rods of unit length.  The first
rod makes an angle \(\theta\) with the downward vertical.  For the second
degree of freedom we have a choice.  We could use a second absolute angle
\(\phi\), or we could use the relative angle \(\alpha\) between the rods.

The relative angle is the more revealing choice.  Temporarily turn gravity
off.  A rigid rotation of the complete apparatus changes the absolute
orientation but not the angle between the rods:
\begin{equation}
  \theta\longrightarrow\theta+\epsilon,
  \qquad
  \alpha\longrightarrow\alpha.
  \label{eq:l6-double-rotation}
\end{equation}
With this choice the rotational symmetry acts on one coordinate only.  This
does not make other coordinate systems wrong.  It makes the symmetry, and
therefore the conserved quantity, easier to see.

\begin{figure}[htbp]
  \centering
  \includegraphics[width=0.76\textwidth]{lecture_06_figure_09.png}
  \caption{The double pendulum with \(\theta\) as the first absolute angle and
  \(\alpha\) as the second rod's angle relative to the first.}
  \lectureframe{00:29:20}
\end{figure}

To obtain the kinetic energy, use Cartesian coordinates as temporary
bookkeeping devices.  With the vertical coordinate measured downward, the
first bob is at
\begin{equation}
  x_1=\sin\theta,
  \qquad
  y_1=\cos\theta,
  \label{eq:l6-bob-one-position}
\end{equation}
and hence
\begin{equation}
  \dot x_1=\dot\theta\cos\theta,
  \qquad
  \dot y_1=-\dot\theta\sin\theta,
  \qquad
  \dot x_1^2+\dot y_1^2=\dot\theta^2.
  \label{eq:l6-bob-one-velocity}
\end{equation}
The second bob is at
\begin{align}
  x_2&=\sin\theta+\sin(\theta+\alpha),\\
  y_2&=\cos\theta+\cos(\theta+\alpha),
  \label{eq:l6-bob-two-position}
\end{align}
so, for example,
\begin{equation}
  \dot x_2=\dot\theta\cos\theta
  +(\dot\theta+\dot\alpha)\cos(\theta+\alpha),
  \label{eq:l6-bob-two-xvelocity}
\end{equation}
with an analogous expression for \(\dot y_2\).

\begin{figure}[htbp]
  \centering
  \includegraphics[width=0.88\textwidth]{lecture_06_figure_10.png}
  \caption{Auxiliary Cartesian coordinates for the two bobs.  They are used
  to compute the velocities, not introduced as additional degrees of freedom.}
  \lectureframe{00:35:25}
\end{figure}

Adding the two kinetic energies gives the consistent result below.  The cross
term follows from
\begin{equation*}
  \cos A\cos B+\sin A\sin B=\cos(A-B).
\end{equation*}
\begin{equation}
  T=\dot\theta^2
    +\frac12(\dot\theta+\dot\alpha)^2
    +\dot\theta(\dot\theta+\dot\alpha)\cos\alpha.
  \label{eq:l6-double-kinetic}
\end{equation}

\begin{classroomqa}
  \audiencequestion{Does this expression already include the kinetic energy
  of both bobs, and has a factor of one half been lost?
  \lecturetimestamp{00:38:06--00:38:27}}
  \lecturerresponse{It includes both.  It comes from
  \(\tfrac12(\dot x_1^2+\dot y_1^2+\dot x_2^2+\dot y_2^2)\).
  Two equal contributions combine to produce the first term
  \(\dot\theta^2\); no factor is missing.}
\end{classroomqa}

Without gravity, \(L=T\) has no explicit \(\theta\).  Thus
\begin{equation}
  \dot p_\theta=\frac{\partial L}{\partial\theta}=0,
  \label{eq:l6-double-angular-conservation}
\end{equation}
and \(p_\theta\) is the total angular momentum.  Notice the precise meaning of
the statement that \(\theta\) does not appear: \(\dot\theta\) certainly
appears, but the undifferentiated coordinate does not.

Restoring gravity breaks the global rotational symmetry.  For the downward
vertical convention above, one convenient choice of zero gives
\begin{equation}
  V=-g(y_1+y_2)
   =-g\left[2\cos\theta+\cos(\theta+\alpha)\right].
  \label{eq:l6-double-potential}
\end{equation}
An additive constant would change nothing.  The important point is that
\(\theta\) now appears explicitly, so its conjugate momentum is no longer
conserved.  Energy remains conserved because the Lagrangian has no explicit
time dependence.

\begin{classroomqa}
  \audiencequestion{Should the board's
  \(\dot\theta-\dot\alpha\) be
  \(\dot\theta+\dot\alpha\)?
  \lecturetimestamp{00:43:00--00:43:54}}
  \lecturerresponse{The orientation used here requires the plus sign.  It is
  the convention in Eq.~\eqref{eq:l6-bob-two-position} and is carried through
  the kinetic energy in Eq.~\eqref{eq:l6-double-kinetic}.  Reversing the
  definition of positive \(\alpha\) changes the corresponding signs, but every
  occurrence must then be changed consistently.  The board mixed the two
  conventions.}
\end{classroomqa}

\begin{figure}[htbp]
  \centering
  \includegraphics[width=0.84\textwidth]{lecture_06_figure_02.png}
  \caption{The no-gravity pendulum calculation and
  \(\dot p_\theta=\partial L/\partial\theta\).  The upper board line preserves
  the sign-inconsistent local notation discussed in the question; the
  normalized formula is Eq.~\eqref{eq:l6-double-kinetic}.}
  \lectureframe{00:45:50}
\end{figure}

Once \(T\) and \(V\) are known, the remaining equations are obtained by
differentiation.  Rod tensions and the forces that enforce fixed lengths have
again disappeared into the choice of coordinates.

\subsection{What conservation does and does not imply}

\begin{classroomqa}
  \audiencequestion{Can one see directly from the equations whether the
  double pendulum is chaotic?
  \lecturetimestamp{00:44:57--00:45:25}}
  \lecturerresponse{Usually not by simple inspection.  As parameters or
  energies vary, a system may pass from regular to chaotic behavior.  The
  subject requires more structure than the raw appearance of the equations,
  and it will be taken up later.}
\end{classroomqa}

\begin{classroomqa}
  \audiencequestion{In a conservative system, is every \(\dot p_i\) equal to
  zero?
  \lecturetimestamp{00:45:25--00:46:24}}
  \lecturerresponse{No.  The general equation is
  \(\dot p_i=\partial L/\partial q_i\).  It vanishes only when \(q_i\) is
  cyclic.  A cyclic coordinate is the especially simple case in which a
  symmetry has been made to shift one coordinate while leaving the others
  fixed.}
\end{classroomqa}

If the only objective is numerical evolution, an inconvenient coordinate
system still works; it merely produces longer equations.  The reward for a
clever system is conceptual: the conserved momentum is visible before the
equations are solved.

\begin{classroomqa}
  \audiencequestion{Without gravity, does a double pendulum behave as simply
  as a single free pendulum?
  \lecturetimestamp{00:47:02--00:47:35}}
  \lecturerresponse{A single free rotor turns uniformly because its one
  angular momentum fixes its motion.  A double pendulum can execute very
  complicated relative motion while preserving the total angular momentum.
  If both angular velocities initially vanish it remains at rest; otherwise
  the two parts interact.}
\end{classroomqa}

\begin{classroomqa}
  \audiencequestion{Are \(\dot\theta\) and \(\dot\alpha\) individually
  constant when gravity is absent?
  \lecturetimestamp{00:47:36--00:52:42}}
  \lecturerresponse{No.  Conservation says \(\dot p_\theta=0\), not
  \(\ddot\theta=0\) or \(\ddot\alpha=0\).  The expression for
  \(p_\theta=\partial L/\partial\dot\theta\) contains both angular velocities
  and depends on \(\alpha\).  The rods transmit torques, allowing the two bobs
  to exchange angular momentum and energy.  This is analogous to two masses
  joined by a spring: total momentum and center-of-mass velocity are constant,
  while the individual velocities oscillate.  The rotational analogue is a
  particular collective combination, not either angle by itself.}
\end{classroomqa}

\begin{classroomqa}
  \audiencequestion{Are the forces that exchange angular momentum carried
  through the rods?
  \lecturetimestamp{00:52:43--00:53:34}}
  \lecturerresponse{Yes.  A force transmitted along a rod need not point
  through the pivot of the other motion, so it can exert a torque.  A torque
  is simply the rotational tendency of a force, and here it produces the
  nonzero individual angular accelerations.}
\end{classroomqa}

\section{The Hamiltonian as a Third Formulation}

The Hamiltonian is energy, but that is only its first significance.  It is
also the basis of a third formulation of mechanics.  Newton's equations use
forces and accelerations.  The principle of stationary action gives the
Euler--Lagrange equations.  Hamiltonian mechanics is equivalent to the latter
for regular systems, but its variables and geometry look different.

Instead of \(q_i\) and \(\dot q_i\), use \(q_i\) and their conjugate momenta
\(p_i\).  A state of a system with \(n\) degrees of freedom requires \(2n\)
numbers in either description: positions and velocities, or positions and
momenta.  The bridge between them is
\begin{equation}
  p_i=\frac{\partial L}{\partial\dot q_i}.
  \label{eq:l6-legendre-map}
\end{equation}
When these relations can be inverted for the velocities
\(\dot q_i(q,p)\), define
\begin{equation}
  H(q,p)=\sum_i p_i\dot q_i-L(q,\dot q).
  \label{eq:l6-hamiltonian-definition}
\end{equation}
The velocities on the right are then eliminated in favor of \(q\) and \(p\).
This invertibility is the regular Legendre-transform assumption.  Singular
Lagrangians require the more general theory of constraints; they are not
handled by this elementary construction.

\subsection{Reversibility before calculus}

Recall a finite collection of abstract states connected by arrows.  A
deterministic rule may still fail to be reversible if two earlier states flow
into one later state.  The future is unique, but the past is not.  Reversing
all the arrows gives the opposite defect: one state branches into several
futures.  In the first picture arrows converge; in the second they diverge.

\begin{figure}[htbp]
  \centering
  \includegraphics[width=0.86\textwidth]{lecture_06_figure_11.png}
  \caption{Abstract state maps used to contrast reversible cycles with
  forbidden convergence and divergence.}
  \lectureframe{00:56:20}
\end{figure}

Hamiltonian flow will supply the continuous version of this statement.  A
cloud of phase-space points may move and deform, but it cannot acquire a sink
that destroys distinctions between states or a source that manufactures
several futures from one state.  The proof is deferred; for now the diagrams
motivate the geometry.

\begin{classroomqa}
  \audiencequestion{Should the abstract state diagram be interpreted as an
  object moving in a gravitational field?
  \lecturetimestamp{00:58:02--00:58:35}}
  \lecturerresponse{No.  The dots are merely abstract states and the arrows
  specify a rule for changing state.  No spatial or gravitational meaning has
  been assigned to the drawing.}
\end{classroomqa}

\begin{classroomqa}
  \audiencequestion{How is a forbidden convergent state map distinguished
  from a black hole?
  \lecturetimestamp{00:58:40--00:59:28}}
  \lecturerresponse{Only with difficulty and subtlety.  It is a fair question,
  but a correct answer requires a careful account of the complete state and
  its quantum-gravitational evolution.  The elementary diagram is not enough
  to settle it.}
\end{classroomqa}

\section{An Oscillator Adapted to Phase Space}

Before deriving the general equations, consider a harmonic oscillator written
in deliberately adapted variables:
\begin{equation}
  L=\frac{1}{2\omega}\dot q^2-\frac{\omega}{2}q^2.
  \label{eq:l6-rescaled-lagrangian}
\end{equation}
Its conjugate momentum is
\begin{equation}
  p=\frac{\partial L}{\partial\dot q}=\frac{\dot q}{\omega},
  \qquad
  \dot q=\omega p.
  \label{eq:l6-rescaled-momentum}
\end{equation}
Substitution in Eq.~\eqref{eq:l6-hamiltonian-definition} gives
\begin{align}
  H&=p\dot q-L\\
   &=\omega p^2-\frac{\omega p^2}{2}+\frac{\omega q^2}{2}
    =\frac{\omega}{2}(p^2+q^2).
  \label{eq:l6-oscillator-hamiltonian}
\end{align}

\begin{figure}[htbp]
  \centering
  \includegraphics[width=0.86\textwidth]{lecture_06_figure_12.png}
  \caption{The Legendre transform of the adapted oscillator and the symmetric
  Hamiltonian \(H=\tfrac{\omega}{2}(p^2+q^2)\).}
  \lectureframe{01:08:20}
\end{figure}

The space of \(q\) and \(p\) is phase space.  One degree of freedom gives a
two-dimensional plane; two degrees of freedom give a four-dimensional phase
space.  At fixed oscillator energy,
\begin{equation}
  p^2+q^2=\frac{2H}{\omega},
  \label{eq:l6-energy-circle}
\end{equation}
so the trajectory remains on a circle.  Energy conservation confines the
motion to one such curve.

\begin{figure}[htbp]
  \centering
  \includegraphics[width=0.82\textwidth]{lecture_06_figure_13.png}
  \caption{Nested constant-energy circles in the oscillator phase plane,
  with \(q\) horizontal and \(p\) vertical.}
  \lectureframe{01:13:20}
\end{figure}

\section{Deriving Hamilton's Equations}

The derivation is an exercise in the differential of a function of several
variables.  For \(F(q,p)\),
\begin{equation}
  \delta F=\sum_i\left(
  \frac{\partial F}{\partial q_i}\delta q_i
  +\frac{\partial F}{\partial p_i}\delta p_i\right).
  \label{eq:l6-total-differential}
\end{equation}
The coefficients can be read in either direction: the coefficient of
\(\delta q_i\) is \(\partial F/\partial q_i\), and the coefficient of
\(\delta p_i\) is \(\partial F/\partial p_i\).

Now vary the defining expression for \(H\):
\begin{align}
  \delta H
  &=\sum_i\left(\dot q_i\,\delta p_i
       +p_i\,\delta\dot q_i\right)-\delta L\\
  &=\sum_i\left(\dot q_i\,\delta p_i
       +p_i\,\delta\dot q_i
       -\frac{\partial L}{\partial q_i}\delta q_i
       -\frac{\partial L}{\partial\dot q_i}\delta\dot q_i\right).
  \label{eq:l6-vary-hamiltonian}
\end{align}
Because \(\partial L/\partial\dot q_i=p_i\), the two
\(\delta\dot q_i\) terms cancel.  The Euler--Lagrange equations give
\(\partial L/\partial q_i=\dot p_i\), leaving
\begin{equation}
  \delta H=\sum_i\left(\dot q_i\,\delta p_i
                     -\dot p_i\,\delta q_i\right).
  \label{eq:l6-reduced-variation}
\end{equation}

\begin{figure}[htbp]
  \centering
  \includegraphics[width=0.86\textwidth]{lecture_06_figure_14.png}
  \caption{The product-rule expansion of \(\delta H\) before the
  \(p_i\,\delta\dot q_i\) terms cancel.}
  \lectureframe{01:18:40}
\end{figure}

Comparing Eq.~\eqref{eq:l6-reduced-variation} with the general differential
in Eq.~\eqref{eq:l6-total-differential}, the independent coefficients give
Hamilton's equations:
\begin{equation}
  \boxed{
  \dot q_i=\frac{\partial H}{\partial p_i},
  \qquad
  \dot p_i=-\frac{\partial H}{\partial q_i}.}
  \label{eq:l6-hamilton-equations}
\end{equation}

\begin{figure}[htbp]
  \centering
  \includegraphics[width=0.78\textwidth]{lecture_06_figure_03.png}
  \caption{The Hamiltonian definition, its variation, and the boxed canonical
  equations.}
  \lectureframe{01:22:18}
\end{figure}

For \(n\) coordinates these are \(2n\) first-order equations.  The
Euler--Lagrange formulation instead supplies \(n\) second-order equations.
The amount of initial data and the physical content are identical.

\section{Oscillator Motion as a Rigid Rotation}

Applying Eq.~\eqref{eq:l6-hamilton-equations} to
Eq.~\eqref{eq:l6-oscillator-hamiltonian} gives
\begin{equation}
  \dot q=\omega p,
  \qquad
  \dot p=-\omega q.
  \label{eq:l6-oscillator-first-order}
\end{equation}
Differentiating the first equation and using the second,
\begin{equation}
  \ddot q=\omega\dot p=-\omega^2q.
  \label{eq:l6-oscillator-second-order}
\end{equation}
Thus the two first-order equations are exactly the usual second-order
oscillator equation, together with the definition of momentum.

In phase space the velocity vector at \((q,p)\) is
\((\dot q,\dot p)=(\omega p,-\omega q)\), tangent to the energy circle.
Every radius rotates with the same angular speed \(\omega\).  The entire plane
may be imagined as a rigid pinwheel or as an incompressible fictitious fluid.
Its points circulate; they do not squeeze into a sink or fly apart from a
source.

\begin{figure}[htbp]
  \centering
  \includegraphics[width=0.86\textwidth]{lecture_06_figure_04.png}
  \caption{A cloud of initial points distributed over nested phase-space
  trajectories.  Hamiltonian evolution moves the cloud rather than collapsing
  it into a preferred state.}
  \lectureframe{01:39:27}
\end{figure}

\begin{figure}[htbp]
  \centering
  \includegraphics[width=0.90\textwidth]{lecture_06_figure_05.png}
  \caption{The oscillator phase portrait beside the equivalent
  Euler--Lagrange equation
  \(\ddot q/\omega=-\omega q\).}
  \lectureframe{01:42:10}
\end{figure}

The next step will be to turn the fluid analogy into a theorem: Hamiltonian
flow has no local convergence or divergence.  That result reconnects the
continuous mechanics with the reversible state maps considered at the
beginning of the course.

\section{Questions That Expose the Coordinate Choice}

\begin{classroomqa}
  \audiencequestion{Is the assumption that velocities can be solved in terms
  of \(q\) and \(p\) a theorem or an empirical property of physical systems?
  \lecturetimestamp{01:29:08--01:29:45}}
  \lecturerresponse{Within this elementary development it is an empirical
  regularity, with a deeper motivation from quantum mechanics.  More
  precisely, Eq.~\eqref{eq:l6-legendre-map} must have a nonsingular Jacobian
  with respect to the velocities.  Modern mechanics also contains important
  singular Lagrangians, especially constrained and gauge systems, but their
  Hamiltonian formulation needs additional constraint machinery.}
\end{classroomqa}

An objection about units explains why the oscillator was written in the
unfamiliar form of Eq.~\eqref{eq:l6-rescaled-lagrangian}.  Begin instead with
the ordinary coordinate \(x\):
\begin{equation}
  L=\frac{m}{2}\dot x^2-\frac{k}{2}x^2,
  \qquad
  \omega=\sqrt{\frac{k}{m}}.
  \label{eq:l6-ordinary-oscillator}
\end{equation}
Choose
\begin{equation}
  q=(km)^{1/4}x.
  \label{eq:l6-oscillator-rescaling}
\end{equation}
Then \(q^2=\sqrt{km}\,x^2\) and
\(\dot q^2=\sqrt{km}\,\dot x^2\), which converts
Eq.~\eqref{eq:l6-ordinary-oscillator} precisely into
Eq.~\eqref{eq:l6-rescaled-lagrangian}.  The new \(q\) is an adapted generalized
coordinate, not a length, and its conjugate \(p\) is not the ordinary
mechanical momentum \(p_x=m\dot x\).

\begin{figure}[htbp]
  \centering
  \includegraphics[width=0.86\textwidth]{lecture_06_figure_15.png}
  \caption{The dimensional rescaling from the usual oscillator coordinate
  \(x\) to the adapted coordinate \(q=(km)^{1/4}x\).}
  \lectureframe{01:36:10}
\end{figure}

\begin{classroomqa}
  \audiencequestion{How can \(\omega p=\dot q\) be dimensionally consistent,
  and why were these variables chosen?
  \lecturetimestamp{01:29:46--01:39:22}}
  \lecturerresponse{The symbol \(q\) is not required to have dimensions of
  length, and its conjugate \(p\) is not required to have the dimensions of
  \(m\dot x\).  The rescaling in Eq.~\eqref{eq:l6-oscillator-rescaling}
  preserves the dimensions of the Lagrangian and makes \(q\) and \(p\) enter
  \(H\) symmetrically.  In the original variables,
  \[
    H=\frac{p_x^2}{2m}+\frac{kx^2}{2},
  \]
  so constant-energy curves in the \((x,p_x)\) plane are ellipses.  In the
  adapted \((q,p)\) plane they are circles.  The rescaling was chosen to make
  that symmetry visible.}
\end{classroomqa}

\begin{classroomqa}
  \audiencequestion{How can \(2n\) first-order Hamilton equations be
  equivalent to \(n\) second-order Lagrange equations?
  \lecturetimestamp{01:39:24--01:42:30}}
  \lecturerresponse{For the oscillator, differentiate
  \(\dot q=\omega p\) and substitute \(\dot p=-\omega q\); the pair becomes
  \(\ddot q=-\omega^2q\).  Starting directly from the Lagrangian gives
  \[
    \frac{d}{dt}\left(\frac{\dot q}{\omega}\right)
      =-\omega q,
  \]
  which is the same second-order equation.  One Hamilton equation can be
  viewed as defining momentum from velocity, while the other evolves that
  momentum.}
\end{classroomqa}

\section{Why Phase Space and Potential Energy Matter}

\begin{classroomqa}
  \audiencequestion{Does the power of the Hamiltonian mean that phase space
  is a better geometrical representation than configuration space?
  \lecturetimestamp{01:42:33--01:44:38}}
  \lecturerresponse{They serve different purposes.  Phase space makes cycles,
  constant-energy surfaces, and the absence of convergence or divergence
  unusually transparent.  Configuration space may be better for other
  questions.  The deeper reason for emphasizing \(H\) is that it becomes a
  central and more intuitive object in quantum mechanics, one of the modern
  subjects toward which this mechanics is directed.}
\end{classroomqa}

\begin{classroomqa}
  \audiencequestion{Should potential energy be regarded as the integral of
  force?
  \lecturetimestamp{01:44:40--01:45:45}}
  \lecturerresponse{For conservative forces that statement is valid, up to a
  constant and with the appropriate sign.  The more useful direction is to
  regard the potential as primary and derive the generalized force from it:
  \[
    F_i=-\frac{\partial V}{\partial q_i}.
  \]
  This is exactly the term that appears in
  \(\dot p_i=\partial L/\partial q_i\) when \(L=T-V\) and the kinetic energy
  has no explicit dependence on \(q_i\).}
\end{classroomqa}

\begin{classroomqa}
  \audiencequestion{Which frontier problems is this path toward modern
  physics meant to illuminate?
  \lecturetimestamp{01:45:46--01:47:47}}
  \lecturerresponse{The path leads toward quantum gravity and cosmology;
  elementary-particle physics, including the Higgs sector and neutrinos;
  relativity; string theory; and quantum field theory.  The enduring
  theoretical problem is to fit these pieces into one coherent structure,
  especially quantum mechanics with gravity and cosmology with particle
  physics.  This was a statement of the frontier as it appeared in 2011, not a
  claim that its solution was already known.}
\end{classroomqa}

The two constrained examples and the oscillator now fit into one picture.
Generalized coordinates remove unnecessary forces.  Symmetry selects
conserved canonical momenta.  The Legendre transform exchanges velocities for
momenta, and Hamilton's equations turn dynamics into a flow on phase space.
For the oscillator that flow is a rigid rotation; for a general system its
geometry will provide the precise continuous meaning of reversible motion.
